% !TEX program = xelatex
\documentclass[11pt,a4paper]{ctexart}

% ==================== 颜色 ====================
\usepackage{xcolor}
\definecolor{maincolor}{HTML}{1A7A6B}       % 主色（标题、顶栏、框边框）
\definecolor{accent}{HTML}{D4A843}          % 强调色（关键词高亮）
\definecolor{boxbg}{HTML}{F0F7F5}           % 问题框背景
\definecolor{borderblue}{HTML}{2E86AB}      % 方法框边框
\definecolor{proofbg}{HTML}{FDF8F0}         % 方法框背景
\definecolor{summarybg}{HTML}{F5F0FF}       % 总结框背景

% ==================== 页面 ====================
\usepackage[top=2.2cm, bottom=2.2cm, left=2.5cm, right=2.5cm, headheight=23.12pt]{geometry}

% ==================== 字体 ====================
\setmainfont{TeX Gyre Pagella}

% ==================== 数学 ====================
\usepackage{amsmath,amssymb,amsfonts,mathtools}

% ==================== 彩色框 ====================
\usepackage{tcolorbox}
\tcbuselibrary{skins,breakable,most}

% ==================== 页眉页脚 ====================
\usepackage{fancyhdr}
\pagestyle{fancy}
\fancyhf{}
\fancyhead[L]{\color{maincolor}\sffamily\bfseries 数学讲义}
\fancyhead[R]{\color{gray}\sffamily \today}
\fancyfoot[C]{\color{gray}\sffamily\thepage}
\renewcommand{\headrulewidth}{1.5pt}
\renewcommand{\headrule}{\hbox to\headwidth{\color{maincolor}\leaders\hrule height \headrulewidth\hfill}}

% ==================== 章节标题样式 ====================
\usepackage{titlesec}
\titleformat{\section}
  {\color{maincolor}\sffamily\LARGE\bfseries}
  {\thesection}{1em}{}
  [\vspace{0.3cm}{\color{maincolor!40}\titlerule[2pt]}]

\titleformat{\subsection}
  {\color{maincolor!80}\sffamily\Large\bfseries}
  {\thesubsection}{1em}{}

% ==================== 表格 ====================
\usepackage{array,booktabs,colortbl}

% ==================== 杂项 ====================
\usepackage{enumitem}
\setlist{nosep, leftmargin=*}
\usepackage[hidelinks]{hyperref}

% ==================== 自定义命令 ====================
\newcommand{\highlight}[1]{\colorbox{accent!20}{\sffamily\bfseries #1}}

% ==================== 自定义环境 ====================

% 1. 问题框 —— 青绿边框 + 浅绿背景
\newtcolorbox[auto counter]{problembox}[1][]{
  enhanced,
  colframe=maincolor,
  colback=boxbg,
  coltitle=white,
  fonttitle=\sffamily\bfseries,
  title={\large 问题},
  attach boxed title to top left={xshift=4mm, yshift=-2mm},
  boxed title style={colback=maincolor, sharp corners},
  sharp corners,
  boxrule=1.2pt,
  left=4mm, right=4mm, top=5mm, bottom=4mm,
  before skip=1.2em, after skip=1.2em,
  #1
}

% 2. 方法/证明框 —— 蓝边框 + 暖白背景，第二个参数是标题文字
\newtcolorbox[auto counter]{methodbox}[2][]{
  enhanced,
  colframe=borderblue,
  colback=proofbg,
  coltitle=white,
  fonttitle=\sffamily\bfseries,
  title={#2},
  attach boxed title to top left={xshift=4mm, yshift=-2mm},
  boxed title style={colback=borderblue, sharp corners},
  sharp corners,
  boxrule=1.2pt,
  left=4mm, right=4mm, top=5mm, bottom=4mm,
  before skip=1.2em, after skip=1.2em,
  #1
}

% 3. 总结框 —— 淡紫背景
\newtcolorbox{summarybox}[1][]{
  enhanced,
  colframe=maincolor!60,
  colback=summarybg,
  coltitle=white,
  fonttitle=\sffamily\bfseries,
  title={\large $\blacktriangleright$ 总结},
  attach boxed title to top left={xshift=4mm, yshift=-2mm},
  boxed title style={colback=maincolor!70, sharp corners},
  sharp corners,
  boxrule=1.2pt,
  left=4mm, right=4mm, top=5mm, bottom=4mm,
  before skip=1.5em, after skip=1.2em,
  #1
}

% ==================== 文档开始 ====================
\begin{document}

% ---- 彩色顶栏标题 ----
\begin{tcolorbox}[
  enhanced,
  colframe=maincolor,
  colback=maincolor,
  sharp corners,
  boxrule=0pt,
  height=2.8cm,
  valign=center,
  before skip=-2.2cm,
  after skip=1.5cm,
  grow to left by=2.5cm,
  grow to right by=2.5cm,
]
  \begin{center}
    {\color{white}\sffamily\bfseries\Huge 完全平方数问题}\\[3pt]
    {\color{white!85}\sffamily\large 求所有正整数 $n$ 使 $n^2 + 2n + 12$ 为完全平方数}
  \end{center}
\end{tcolorbox}

\vspace{0.5cm}

\begin{flushright}
  {\color{gray}\sffamily \today}
\end{flushright}

% ==================== 正文 ====================

\section{问题陈述}

\begin{problembox}
求所有正整数 $n$，使得表达式
\[
  n^2 + 2n + 12
\]
为完全平方数（即等于某个整数的平方）。
\end{problembox}

本题是一道典型的初等数论题，考察完全平方数的构造与不等式放缩技巧。下面给出两种不同思路的解法。

\section{解法一：配方与因式分解法}

\begin{methodbox}{核心思路}
先将二次表达式配方，引入新变量 \highlight{$m$} 表示目标平方数，再通过移项、平方差公式化为两因式相乘等于一个\highlight{质数}的形式，利用质数的因子唯一性建立方程组求解。
\end{methodbox}

\subsection*{第一步：设未知平方数}

假设原式等于某个正整数 $m$ 的平方：
\[
  n^2 + 2n + 12 = m^2
\]

\subsection*{第二步：配方}

对等式左边的前两项进行配方：
\[
  (n^2 + 2n + 1) + 11 = m^2
\]
即
\[
  (n+1)^2 + 11 = m^2
\]

\subsection*{第三步：移项并因式分解}

将 $(n+1)^2$ 移到右边，利用平方差公式：
\begin{align*}
  m^2 - (n+1)^2 &= 11 \\[4pt]
  (m - n - 1)(m + n + 1) &= 11
\end{align*}

我们得到了一个\highlight{关键等式}：两个正整数因式的乘积等于 $11$。

\subsection*{第四步：利用质数性质}

因为 $11$ 是\highlight{质数}，且 $m,n$ 均为正整数，两个因式都是正整数。质数 $11$ 的正因子只有 $1$ 和 $11$。

又由于 $m + n + 1 > m - n - 1$，所以必有：
\[
  \begin{cases}
    m - n - 1 = 1 \\[4pt]
    m + n + 1 = 11
  \end{cases}
\]

\subsection*{第五步：求解 $n$}

两式相减消去 $m$：
\begin{align*}
  (m + n + 1) - (m - n - 1) &= 11 - 1 \\
  2n + 2 &= 10 \\
  n &= 4
\end{align*}

代入验证：$n=4$ 时，$n^2 + 2n + 12 = 16 + 8 + 12 = 36 = 6^2$，确为完全平方数。

% ==================== 新页 ====================
\newpage

\section{解法二：放缩与夹逼法}

\begin{methodbox}{核心思路}
不引入额外变量，直接通过不等式放缩，将原式\highlight{夹在}两个相邻（次邻）的完全平方数之间。若原式本身是完全平方数，则它必须等于两界之间唯一可能的那个平方数，从而直接解出 $n$。
\end{methodbox}

\subsection*{第一步：寻找下界}

比较 $n^2 + 2n + 12$ 与 $(n+1)^2$：
\begin{align*}
  (n+1)^2 &= n^2 + 2n + 1
\end{align*}

显然
\[
  n^2 + 2n + 12 > n^2 + 2n + 1 = (n+1)^2
\]

这里没有使用不等式方向——差值 $11 > 0$，对任意正整数 $n$ 恒成立。于是我们得到严格的\highlight{下界}：
\[
  \boxed{n^2 + 2n + 12 > (n+1)^2}
\]

\subsection*{第二步：寻找上界}

再比较 $n^2 + 2n + 12$ 与 $(n+3)^2$：
\begin{align*}
  (n+3)^2 &= n^2 + 6n + 9
\end{align*}

计算差值：
\begin{align*}
  (n+3)^2 - (n^2 + 2n + 12) &= (n^2 + 6n + 9) - (n^2 + 2n + 12) \\
  &= 4n - 3
\end{align*}

当 $n \geqslant 1$ 时，$4n - 3 \geqslant 1 > 0$，恒成立。因此有\highlight{上界}：
\[
  \boxed{n^2 + 2n + 12 < (n+3)^2}
\]

\subsection*{第三步：夹逼}

综合上下界：
\[
  (n+1)^2 < n^2 + 2n + 12 < (n+3)^2
\]

原式被严格夹在两个平方数之间。在正整数序列中，$(n+1)^2$ 与 $(n+3)^2$ 之间只有\highlight{一个}完全平方数——$(n+2)^2$。

\subsection*{第四步：唯一解}

因此，若原式本身是完全平方数，则它\highlight{必须}等于 $(n+2)^2$：
\[
  n^2 + 2n + 12 = (n+2)^2
\]

展开右边并化简：
\begin{align*}
  n^2 + 2n + 12 &= n^2 + 4n + 4 \\
  2n + 12 &= 4n + 4 \\
  2n &= 8 \\
  n &= 4
\end{align*}

同样得到 $n=4$，与解法一结果一致。

\section{总结}

\begin{summarybox}

\subsection*{方法对比}

\begin{tabular}{@{}>{\sffamily\bfseries}p{2.6cm} p{5.5cm} p{5.5cm}@{}}
\toprule
& \color{maincolor} \textbf{配方与因式分解法} & \color{borderblue} \textbf{放缩与夹逼法} \\
\midrule
\textbf{切入点}   & 设未知平方数 $m$，构造等式 & 直接比较原式与邻近平方数 \\
\textbf{核心工具} & 配方法、平方差公式、质数性质 & 不等式放缩、平方数间隔 \\
\textbf{关键步骤} & 化为 $(m-n-1)(m+n+1)=11$ & 夹在 $(n+1)^2$ 与 $(n+3)^2$ 之间 \\
\textbf{技巧要点} & 需要注意到 $11$ 是质数 & 上界选择 $(n+3)^2$ 需试凑 \\
\bottomrule
\end{tabular}

\vspace{1em}

\subsection*{结论}

两种方法殊途同归，原式 $n^2 + 2n + 12$ 为完全平方数当且仅当：

\begin{center}
\begin{tcolorbox}[
  enhanced, colframe=accent!80, colback=accent!10,
  sharp corners, boxrule=1.5pt, width=6cm, center,
]
  {\color{maincolor}\sffamily\bfseries\LARGE $n = 4$}
\end{tcolorbox}
\end{center}

验证：$4^2 + 2 \times 4 + 12 = 16 + 8 + 12 = 36 = 6^2$。

\end{summarybox}

\end{document}
